
\subsection{Formalisation}
	GT-polytopes are represented as a row of integers, with a colon separating 
	the integers into two groups. 
	The integers before the colon represents the top row, and the integers after 
	the colon represents the difference in row sum between the rows. 
	
	In order to create a correct GT-pattern, if the number of differences between row sums 
	exceeds the number of rows created by the top row, zeros are added to the top row. 
	\begin{example}
		Consider the GT-polytope 
		$3\text{ }2\text{ }1 : 2\text{ }1\text{ }1\text{ }1\text{ }1$. 
		In this case, we have 5 differences of row sum and thus 5 rows are required.  
		The top row then becomes $\lambda = [3\text{ }2\text{ }1\text{ }0\text{ }0]$. 
		
		The GT-pattern becomes 
		\begin{center}
		\begin{tikzpicture}
			\node at (-4, 4) {$3$};
			\node at (-2, 4) {$2$};
			\node at (0, 4) {$1$};
			\node at (2, 4) {$0$};
			\node at (4, 4) {$0$};
			
			\node at (-3, 3) {$x_{41}$};
			\node at (-1, 3) {$x_{42}$};
			\node at (1, 3) {$x_{43}$};
			\node at (3, 3) {$x_{44}$};
			
			\node at (-2, 2) {$x_{31}$};
			\node at (0, 2) {$x_{32}$};
			\node at (2, 2) {$x_{33}$};
			
			\node at (-1, 1) {$x_{21}$};
			\node at (1, 1) {$x_{22}$};
			
			\node at (0,0) {$x_{11}.$};
			
			\node at (-5, 3.5) {$2$};
			\node at (-4, 2.5) {$1$};
			\node at (-3, 1.5) {$1$};
			\node at (-2, 0.5) {$1$};
			\node at (-1, -0.5) {$1$};
		\end{tikzpicture}
		\end{center}
		In the picture the differences between the row sums are indicated to the left 
		between the rows. 
		For the bottom row, the only variable $x_{11} = 1$ by necessity, since 
		the difference between no row $(=0)$ and the only variable in the bottom row 
		is $1$. 
		
		Continuing upwards we see that the sum for the fourth row from the bottom must be 
		$x_{41} + x_{42} + x_{43} + x_{44} = 4$. The difference between this row and 
		the top row is $2$ and the sum of the top row is indeed $3 + 2 + 1 + 0 + 0 = 6$.
	\end{example}
	
	Some patterns can be ignored. We motivate these:
	
	\begin{enumerate}
		\item Two patterns with a fixed top row $\lambda$ where the summation of the 
		different row sums give the same Ehrhart polynomial.
		\item If the row sum on the left of the colon is not the same as the sum on 
		the right side, we do not get a correct Schur function.
		\item If the difference between two row sums are larger than the greatest 
		integer of the top row in the GT-pattern, i.e. if a number on the right side of a 
		colon is greater than any of the numbers on the left side, we do not get a correct 
		Schur function.
	\end{enumerate}
	Motivation:
	\begin{enumerate}
		\item In Part \ref{gt-polytopes} it was mentioned that the symmetric properties 
		of Schur functions means that for a fixed top row $\lambda$, 
		every GT-polytope where the summation of the different row sums give 
		the same Ehrhart polynomial. This means that only one permutation of the row sum 
		differences need to be checked, as all other permutations will give the same 
		Ehrart polynomial.
		
		\item 
		If they do not add up to the same sum, there are two cases: 
		(a) either the sum on the right side is smaller than on the left or 
		(b) the sum on the right side is larger than on the left. 
		By looking at the bijection between GT-patterns and SSYT's it is easy to see 
		why neither of these cases works out.
		\begin{enumerate}
			\item
			In this case, we would have more boxes in the SSYT than numbers to 
			fill them with. 
			\item In this case, we would have more numbers than boxes to fill them with. 
		\end{enumerate}
		Neither case allows for a correct SSYT. 
		
		\item Once again looking at SSYT's, this is because we would have a greater 
		number of integers of the same size than boxes, and in the downwards direction 
		for a correct SSYT the integers are stricly increasing. Thus, 
		any combination of integers where one or more integers on the right side of 
		the colon was larger than any of the ones on the left side were excluded.
	\end{enumerate} 
	
	
	\begin{example}
		The GT-polytopes 
		$3\text{ }2\text{ }1 : 2\text{ }1\text{ }1\text{ }1\text{ }1$ and 
		$3\text{ }2\text{ }1 : 1\text{ }1\text{ }2\text{ }1\text{ }1$
		give the same Ehrhart polynomial as they both have the top row 
		$\lambda = [3\text{ }2\text{ }1\text{ }0\text{ }0]$ and the row sum differences 
		are just permutations of each other. 
		
		This is motivated by the fact that Schur functions are symmetric.
	\end{example}
	
	\bigskip
	
	The GT-polytopes were generated by choosing an integer and computing all 
	unordered partitions of that integer. These partitions became the top row 
	and was ordered in descending order. Each of these partitions was then paired 
	with all partitions having the greatest number equal to or less than the greatest 
	number of the top row partitions, and these second partitions became the row sum differences. 
	
	\bigskip
	
	Note that no polytopes where the row sum differences where only ones where the sum was ten 
	were not calculated. 
	Those had the most computational power needed to be calculated, and it was not possible with 
	my computer. 
	
\subsection{Realization}
		
	The GT-polytopes were generated using Python\footnote{All code and 
	data can be found at 
	\href{https://github.com/oskarhogman/kandidatprojekt-matte/tree/main/Code}{the Github repository for the project.}}.
	Following the rules above, patterns for an integer $n$ where the sum of the top row 
	$\lambda$ and the sum of the row sum differences both add up to $n$ were created. 
	
	Following this, matrices were generated of each polytope in a format accepted by 
	Julia. The Julia code then calculated the $h^*$-polynomial for each of 
	the GT-polytopes.
	 
	Python was again used to get the Ehrhart polynomial from the $h^*$-polynomial.
	